\documentclass[10pt,a4paper]{amsart}
\usepackage{amsmath,amssymb,amsthm,graphicx}
\RequirePackage{amssymb,amsmath,amsthm,bm}
\RequirePackage{thmtools}
\RequirePackage[capitalize]{cleveref}
\providecommand{\tightlist}{%
  \setlength{\itemsep}{0pt}\setlength{\parskip}{0pt}}

%%%%% varphi instead of phi
\renewcommand{\phi}{\varphi}

\mathchardef\mhyph="2D 
\providecommand{\A}{\mathcal{A}}
\providecommand{\C}{\mathcal{C}}
\providecommand{\L}{\mathcal{L}}
\providecommand{\B}{\mathcal{B}}
\providecommand{\M}{\mathcal{M}}
\renewcommand{\S}{\mathcal{S}}
\providecommand{\ol}[1]{\overline{#1}}
\providecommand{\mbf}[1]{\mathbf{#1}}
\providecommand{\mf}[1]{\mathfrak{#1}}
\providecommand{\mrm}[1]{\mathrm{#1}}
\providecommand{\mc}[1]{\mathcal{#1}}
\providecommand{\ra}{\rightarrow}
\providecommand{\Ra}{\Rightarrow}
\providecommand{\La}{\Leftarrow}
\providecommand{\lr}{\leftrightarrow}
\providecommand{\hf}{\mathbb{HF}}
\providecommand{\spec}{DgSp}
\providecommand{\bSigma}{\bm\Sigma}
\providecommand{\bDelta}{\bm\Delta}
\providecommand{\bPi}{\bm\Pi}
\providecommand{\Pinf}[1]{\Pi^{\mathrm{in}}_{#1}}
\providecommand{\Sinf}[1]{\Sigma^{\mathrm{in}}_{#1}}
\providecommand{\dSinf}[1]{d\mhyph\Sigma^{\mathrm{in}}_{#1}}
\providecommand{\Dinf}[1]{\Delta^{\mathrm{in}}_{#1}}
\providecommand{\bp}{\bar p}
\providecommand{\bq}{\bar q}
\providecommand{\ba}{\bar a}
\providecommand{\bb}{\bar b}
\providecommand{\bc}{\bar c}
\providecommand{\bd}{\bar d}
\providecommand{\bx}{\bar x}
\providecommand{\by}{\bar y}
\providecommand{\bz}{\bar z}
\providecommand{\LO}{\mathrm{LO}}
\declaretheorem[name={Theorem},numberwithin=section]{theorem}
\declaretheorem[name={Proposition}, sibling=theorem]{proposition}
\declaretheorem[name={Lemma}, sibling=theorem]{lemma}
\declaretheorem[name={Fact}, sibling=theorem]{fact}
\declaretheorem[name={Claim}, numberwithin=theorem]{claim}
\declaretheorem[name={Definition},numberwithin=section,sibling=theorem, style=definition]{definition}
\declaretheorem[name={Remark}, numbered=no,style=definition]{remark}
\declaretheorem[name={Corollary}, sibling=theorem]{corollary}
\declaretheorem[name={Question},sibling=theorem]{question}
\newcommand{\T}{{\mathcal{T}}}
\newcommand{\SR}{SR}
\newcommand{\pSR}{pSR}
\newcommand{\Q}{\mathbb{Q}}
\newcommand{\Z}{\mathbb{Z}}
\newcommand{\K}{\mathcal{K}}
\newcommand{\N}{\mathcal{N}}
\renewcommand{\L}{\mathcal{L}}
\providecommand{\Ginf}[1]{\Gamma^{\mathrm{in}}_{#1}}
\title{Countable linear orders of exact Sigma\_(lambda+2) Scott-sentence complexity at every nonzero limit lambda}
\author{David Gonzalez, Dino Rossegger}
\date{arxiv:2305.07126v1; CC BY 4.0}
\usepackage[UTF8,fontset=fandol]{ctex}
\begin{document}
\maketitle
\noindent\textit{Editorial note.} All ordinals are countable; lambda is a nonzero limit and delta\_n is a cofinal fundamental sequence.\par\medskip
\noindent\textit{Editorial note.} The summation variable in1403/1435 is read consistently with the explicitly defined delta\_n fundamental sequence, and the bounded tail in1437 uses delta\_n. Original source notation is preserved.\par\medskip
\noindent\textit{Editorial note.} This construction is not a sum of two shuffle sums. The paper explicitly distinguishes it at1425-1426 and1520-1522 from the already prepared lambda+3 construction; its approximation-sequence and distinguished-cut argument is new work.\par\medskip
\noindent\textit{Editorial note.} The public Z-power back-and-forth lemma141394 is complete support, not a newly counted target.\par\medskip
\noindent\textit{Editorial note.} 2026-10-11材料收口：复用既有H2、来源许可和近邻判断；核对原命题与作者证明起止、所指上下文及装配范围齐备。未重证作者结论；沿用既有非重复依据。\par\medskip
\noindent\textit{Editorial scope.} One original exact result. Native construction and load-bearing local dependencies preserved; named classical Scott-analysis inputs explicit.\par\medskip
\noindent\textit{Allowed standard prerequisites (editorial).}\begin{itemize}
\item Standard infinitary Scott analysis. Orbit-definability/Scott-sentence complexity criteria, interval back-and-forth, and the equality of parameterized rank with the maximum ranks of intervals separated by named points.
\item Classical scattered-order rank bounds. For a scattered order of bounded Hausdorff rank, the standard Alvir-Rossegger/Ash Scott bounds; a parameter-definable interval cannot exceed the ambient Scott rank. Source-specific Z-power comparison and marker construction are provided in full.
\end{itemize}\medskip
In this article we give a comprehensive Scott analysis of the class of linear
orderings using three definitions of rank. Given a structure $\A$ the
\emph{(parameterless) Scott rank} $\SR(\A)$, is the least $\alpha$, such that
the automorphism orbits of all tuples in $\A$ are $\Sinf{\alpha}$ definable; the
\emph{parameterized Scott rank} $\pSR(\A)$, is the least ordinal $\alpha$ such
that there is a parameter $\bp\in A^{<\omega}$ such that $SR((\A,\bp))=\alpha$;
and the \emph{Scott sentence complexity} of $\A$, $SSC(\A)$ is the least of the
complexities $\Pinf{\alpha}$, $\Sinf{\alpha}$, $\dSinf{\alpha}$ such that $\A$
has a Scott sentence of one of these complexities. While the latter is not an

\subsection{Preliminaries}
We assume that the reader is familiar with the basic techniques of Scott
analysis and with linear orderings in general. For background we refer to
Montalb\'an's upcoming book~\cite{montalban2021} and the standard reference by
Ash and Knight~\cite{ash2000}. Let us highlight the following lemma
which we will use as basic tools in our proofs without reference.
\begin{lemma}[{\cite[Lemma 15.8]{ash2000}}]\label{lem:bfandpartitions}  Suppose $\A,\B\in\LO$. Then $\A\leq_1 \B$ if and only if $\A$ is infinite or at least as large as $\B$. For $\alpha>1$, $\A\leq_\alpha \B$ if and only if for every $1\leq\beta<\alpha$ and every partition of $\B$ into intervals $\B_0,\dots,\B_n$, with endpoints in $\B$, there is a partition of $\A$ into intervals $\A_1,\dots,\A_n$ with endpoints in $\A$, such that $\B_i\leq_\beta \A_i$.
\end{lemma}
Notice that this lemma particularly implies that $(\A,\ba)\leq_\alpha (\B,\bb)$
if and only if $(-\infty, a_0)\leq_\alpha (-\infty, b_0)$,
$(a_{|\ba-1|},\infty)\leq_\alpha (b_{|\bb-1|},\infty)$ and for all $0\leq
i<|\ba-1|$ $(a_i,a_{i+1})\leq_\alpha (b_i,b_{i+1})$~\cite[Lemma 15.7]{ash2000}.

In the proof of the below lemma we use $\zeta_\alpha$ to denote the unique up to
isomorphism initial segment of $\Z^\alpha$. 
It is easy to check that $\zeta_\alpha^*$ is the unique end segment of  $\Z^\alpha$ and that $\zeta_{\alpha+1}^*=\zeta_\alpha+\Z^\alpha\cdot\omega$. 

\begin{theorem}\label{thm:Slim+2}
For any limit ordinal $\lambda,$ there is a linear order of Scott sentence
complexity $\Sinf{\lambda+2}.$
\end{theorem}

\begin{lemma}\label{lem:timesZ}
For the sake of organization, consider the following ordinal indexed propositions.
\begin{enumerate}\tightlist
\item[($A_\alpha$)] For any $\K$ and $\L$, $\Z^\alpha\cdot \K
    \leq_{2\alpha}\Z^\alpha\cdot \L$.
\item[($B_\alpha$)] For any $\K$ and $\L$ with $|\K|\geq|L|$, $\Z^\alpha\cdot \K
    \leq_{2\alpha+1}\Z^\alpha\cdot \L$.
\item[($C_\alpha$)] For any $\K$ and any $\L$ without a last element,
    $\Z^\alpha\cdot(\omega+\K)\leq_{2\alpha+2} \Z^\alpha\cdot(\omega+\L)$.
\end{enumerate}
For any countable ordinal $\alpha$, $A_\alpha$, $B_\alpha$ and $C_\alpha$ are true.
\end{lemma}

\begin{proof}
    The proof is by induction with the statements $A_0$ and $B_0$ trivially true.
    We start by showing that $B_\alpha$ implies $A_{\alpha+1}$. 

    Select a tuple $a_0\cdots a_n$ in $\Z^{\alpha+1}\cdot \L=\Z^\alpha\cdot\Z\cdot
    \L$. We will find a tuple $b_0\dots b_n$ in $\Z^{\alpha+1}\cdot
    \K=\Z^\alpha\cdot\Z\cdot \K$ such that $(\Z^{\alpha+1}\cdot \L,\ba
    )\leq_{2\alpha+1} (\Z^{\alpha+1}\cdot \K,\bb)$. Pick a point $b_0$ arbitrarily.
    Given $b_j$, we pick $b_{j+1}$ as follows. If $a_j$ and $a_{j+1}$ are in the
    same copy of $\Z^{\alpha+1}$, we can pick $b_{j+1}$ such that
    $(a_j,a_{j+1})\cong(b_j,b_{j+1})$. If not, we pick $b_{j+1}$
    arbitrarily in the copy of $\Z^{\alpha}$ that is the successor of the successor
    of the copy of $b_j$. In other words, we pick $b_{j+1}$ with  $(b_j,b_{j+1})\cong \zeta^*_\alpha +\Z^\alpha+\zeta_\alpha$.
    
To show that our choice of $\bb$ is correct, we have to confirm the following
relations between non-isomorphic intervals.
\begin{enumerate}\tightlist
\item $(b_j,b_{j+1})\cong \zeta^*_\alpha +\Z^\alpha+\zeta_\alpha \geq_{2\alpha+1} \zeta^*_\alpha
    +\Z^\alpha\cdot \L_1+\zeta_\alpha$, where $\L_1$ is an infinite order,
\item $(-\infty,b_0)\cong \Z^\alpha\cdot \K_1+\zeta_\alpha \geq_{2\alpha+1}\Z^\alpha\cdot
    \L_1+\zeta_\alpha$, where $\L_1$ and $\K_1$ are some initial segments of the $\Z^\alpha$, and therefore infinite, or
\item $(b_{n},\infty)\cong\zeta^*_\alpha +\Z^\alpha\cdot \L_1 \geq_{2\alpha+1} \zeta^*_\alpha
    +\Z^\alpha\cdot \K_1$, where $\L_1$ and $\K_1$ are some final segments of the $\Z^\alpha$, and therefore infinite.
\end{enumerate} 
It follows immediately from $B_\alpha$ that all of these back-and-forth relations hold.

We now show that $B_\alpha$ implies $C_\alpha$. So assume that $\K$ and $\L$ have
no last element and that $B_\alpha$ holds. 
Select a tuple $a_0\cdots  a_n$ in $\Z^\alpha\cdot(\omega+\L)$.
Without loss of generality, $a_0$ is in the first copy of $\Z^\alpha$. We will
find a tuple $b_0\dots b_n$ in $\Z^\alpha\cdot(\omega+\K)$ such that
$(\Z^\alpha\cdot (\omega+\L),\ba)\leq_{2\alpha+1}
(\Z^\alpha\cdot(\omega+\K),\bb)$. For the $a_j$ in the initial $\Z^\alpha\cdot\omega$, the $b_j$ are picked isomorphically. 
We now define the rest of the $b_j$ by induction so that all of them are in the initial $\Z^\alpha\cdot\omega$.
If $a_j$ is in the same $\Z^\alpha$ block as $a_{j+1}$ or $a_{j+1}$ is in the
succesor of $a_j$'s block, define $b_{j+1}$ so that
$(b_j,b_{j+1})\cong (a_j,a_{j+1})$.
Otherwise, define $b_{j+1}$ as being some element in the successor of the
successor of the $\Z^\alpha$ block of $b_j$. 
Not including isomorphic intervals, we only need to check the following cases. 
\begin{enumerate}
    \item $(b_j,b_{j+1})\cong\zeta^*_\alpha +\Z^\alpha+\zeta_\alpha \geq_{2\alpha+1} \zeta^*_\alpha
    +\Z^\alpha\cdot \L_1+\zeta_\alpha$ where $\L_1$ is infinite, 
\item $(b_n,\infty)\cong\zeta^*_\alpha +\Z^\alpha\cdot (\omega+\K) \geq_{2\alpha+1}
    \zeta^*_\alpha +\Z^\alpha\cdot \L_1$, where $\L_1$ is some final segment of
    $\omega+\L$, and thus infinite.
\end{enumerate} 
All of these relations follow from $B_\alpha$, as desired.

To complete the proof of the successor step of the induction, we show that
$C_\alpha$ and $A_{\alpha+1}$ imply $B_{\alpha+1}$.

Select a tuple  $a_0\cdots a_n$ in $\Z^{\alpha+1}\cdot \L$. These points belong
to at most $n$ distinct copies of $\Z^{\alpha+1}$. Because $|\K|\geq| \L|$, we
can match each chosen copy of $\Z^{\alpha+1}$ to one in $\Z^{\alpha+1}\cdot \K$.
Furthermore, if there is a copy of $\Z^{\alpha+1}$ between two of the chosen
copies in $\Z^{\alpha+1}\cdot \L$, we can assure that the corresponding copies
in $\Z^{\alpha+1}\cdot \K$ are also not successors. Also, if $a_n$ is not in the
last copy of $\Z^{\alpha+1}$, we can assure that neither is $b_n$ and, similarly
for $a_0$ and $b_0$ in the first copy. In particular, not including the
isomorphic intervals, we only need to check the following cases.
\begin{enumerate}
    \item $(b_j,b_{j+1})\cong \zeta^*_{\alpha+1} +\Z^{\alpha+1}\cdot
        \K_1+\zeta_{\alpha+1}\geq_{2\alpha+2}\zeta^*_{\alpha+1}
        +\Z^{\alpha+1}\cdot \L_1+\zeta_{\alpha+1}$ for some orders $\K_1$ and $\L_1$,
    \item $(b_n,\infty)\cong\zeta^*_{\alpha+1} +\Z^{\alpha+1}\cdot
        \K_1\geq_{2\alpha+2}\zeta^*_{\alpha+1} +\Z^{\alpha+1}\cdot \L_1$ for
        some orders $\K_1$ and $\L_1$,
\item $(-\infty,b_0)\cong\Z^{\alpha+1}\cdot
    \K_1+\zeta_{\alpha+1}\geq_{2\alpha+2}\Z^{\alpha+1}\cdot
    \L_1+\zeta_{\alpha+1}$ for some orders $\K_1$ and $\L_1$,
\item $(-\infty,b_0)\cong\zeta_{\alpha+1}\geq_{2\alpha+2} \Z^{\alpha+1}\cdot
    \L_1+\zeta_{\alpha+1}$ for some order $\L_1$,
\item $(b_n,\infty)\cong\zeta^*_{\alpha+1} +\Z^{\alpha+1}\cdot
    \K_1\geq_{2\alpha+2} \zeta^*_{\alpha+1} $ for some order $\K_1$.
\end{enumerate}

The first three cases are handled immediately by $A_{\alpha+1}$. By symmetry it
is sufficient to show case (5). However, this is the same as showing that 
$$\zeta_{\alpha}^* + \Z^\alpha\cdot\omega \leq_{2\alpha+2}\zeta_{\alpha}^* +
\Z^{\alpha}\cdot (\omega+\Z\cdot \K_1).$$
This follows directly from $C_\alpha$.




At last we consider the remaining limit levels. We already have seen that
$B_\lambda$ implies $C_\lambda$. The statement $A_\lambda$ follows immediately from $A_\gamma$ for $\gamma<\lambda$. 
A bit more nuanced is the issue of $B_\lambda$. It follows from $A_\lambda$ and $C_\gamma$ for $\gamma<\lambda$.
Analogous to the proof in the successor case, we analyze the possibilities:
\begin{enumerate}
\item $\zeta^*_{\lambda} +\Z^{\lambda}\cdot
    \K_1+\zeta_{\lambda}\geq_\lambda \zeta^*_{\lambda} +\Z^{\lambda}\cdot \L_1+\zeta_{\lambda} $ for some orders $\L_1$ and $\K_1$,
\item $\zeta^*_{\lambda} +\Z^{\lambda}\cdot \K_1\geq_\lambda \zeta^*_{\lambda} +\Z^{\lambda}\cdot \L_1 $ for some orders $\L_1$ and $\K_1$,
\item $\Z^{\lambda}\cdot \K_1+\zeta_{\lambda}\geq_\lambda \Z^{\lambda}\cdot \L_1+\zeta_{\lambda} $ for some orders $\L_1$ and $\K_1$,
\item $\Z^{\lambda}\cdot \K_1+\zeta_{\lambda}\geq_\lambda \zeta_{\lambda} $ for some order $\K_1$,
\item $\zeta^*_{\lambda} +\Z^{\lambda}\cdot \K_1\geq_\lambda\zeta^*_{\lambda} $ for some order $\K_1$.
\end{enumerate}

The first three cases are handled immediately by $A_{\lambda}$. By symmetry it is clear that it is enough to show case 5. However, this is the same as showing that for all $\gamma<\lambda$,
$$\zeta_{\gamma}^* + \Z^\gamma\cdot(\omega +  \Z\cdot\zeta_{\lambda}^*)
\leq_{\gamma}\zeta_{\gamma}^* + \Z^{\gamma}\cdot
(\Z\cdot\zeta_{\lambda}^*+\Z^\lambda\cdot \K_1),$$
which follows from the $C_\gamma$ for $\gamma<\lambda$.

Therefore, for any countable ordinal $\alpha$, $A_\alpha$, $B_\alpha$ and $C_\alpha$ all hold.
\end{proof}

\begin{proposition}\label{prop:scatteredlimit}
For any limit ordinal $\lambda,$ there is a scattered linear order of Scott sentence complexity $\Pinf{\lambda}.$
\end{proposition}

\begin{proof}
Let $(\delta_n)_{n\in\omega}$ be a fundamental sequence for $\lambda$. We show that
\[\L_\lambda:=\sum_{i\in\omega}i+\Z^{\delta_n}\]
is of the desired complexity. 

First note that it indeed does have a $\Pinf{\lambda}$ Scott sentence. The sentence states the following:
\begin{enumerate}\tightlist
    \item There is exactly one 1-block isomorphic to each natural number. 
    \item These 1-blocks are ordered like the natural numbers.
    \item The order between $n$ and $n+1$ is isomorphic to $\Z^{\delta^n}.$
\end{enumerate}
This is $\Pinf{\lambda}$ because each of the items is $\Pinf{<\lambda}$. 
The first two statements can be expressed using finitely many alternations of quantifiers. 
The final statement can be expressed by the Scott sentence of $\Z^{\delta_n}$
relativized to the specific interval, this sentence has Scott sentence
complexity below $\lambda$ by~\cite{AR20}, as the Hausdorff rank of the interval
is below $\lambda$. 

By Lemma~\ref{lem:timesZ}, the $\Z^{\delta_n}$ have unbounded Scott rank below $\lambda$. 
As these are all $\Dinf0$-definable over parameters in $\L_\lambda$,
by~\cite[Lemma 4.3]{MonIntermediate}, we have that for each $n$, $\SR(\L_\lambda)\geq\delta_n,$ so $\Pinf{\lambda}$ is indeed the optimal Scott sentence complexity.
\end{proof}

\begin{proof}
    Let $(\delta_n)_{n\in\omega}$ be a
    fundamental sequence for $\lambda$ and let 
\[\L_\lambda:=\sum_{i\in\omega}i+\Z^{\delta_n}.\]
Furthermore, define for every $n$, the linear orderings:
\[\L_{\lambda,n}:=\sum_{i< n} i+\Z^{\delta_i}+\sum_{i\geq n} i+\Z^{\delta^n}.\]
Note that by Lemma \ref{lem:timesZ}, $\L_\lambda$ is approximated by the
sequence of these orderings in the sense that $\L_\lambda\equiv_{2\delta_n}\L_{\lambda,n}.$

Let $\B=\L_\lambda\cdot(1+\Q)$ and $\A$ be a copy of $\Q$ where an unbounded,
increasing sequence of points $c_n$ is replaced by $\L_{\lambda,n}$ and each
other point is replaced by $\L_\lambda$. We now claim that $\K=\A+\B$ is the
desired order. Note that over the parameter of the initial element of $\B$,
$\SR_p(\K)=\max(\SR(\A),\SR(\B))$ (this equality follows from~\cite[Lemma
11]{GM23}). We now demonstrate that this quantity is exactly $\lambda$.

To define the automorphism orbits of elements within $\B$ we have to distinguish
elements in the initial copy of $\L_\lambda$. For an element $b\in \B$, let $n$
be the size of the finite block to the left of $b$. We can write a formula of
finite rank saying that there is precisely one finite block of size $n$
to the left of $b$ and thus use this type of formulas to distinguish elements in
the first $\L_\lambda$ copy. If $b$ is in a finite block we conjunct this
formula or its negation with its position in this finite block, if $b$ is in a
copy of $\Z^{\delta_n}$, then we relativize the formula defining its
automorphism orbit in this ordering to the interval restricted by block of size
$n$ to the left and a block of size $n+1$ to the right, just as in the proof of
\cref{prop:scatteredlimit}. Conjuncting this with the formula deciding whether
$b$ is in the initial copy of $\L_\lambda$ yields the defining formula for its
automorphism orbit. To extend this definition to tuples we take the conjunction
of the formulas defining orbits for elements and a formula specifying the order
of the tuple.

For $\A$ the definitions are a little more subtle, but generally follow the same
plan. In particular, if the point in $\Q$ that $x$ lies in is between $c_n$ and
$c_{n+1}$, then the following describes the automorphism orbit of $x$:

\begin{enumerate}
    \item Let $z$ be the greatest element below $x$ with no successor or predecessor.

    \item Let $z_n$ be the initial element with no successor or predecessor
        within $c_n$; i.e. say that it is the greatest element with no successor
        or predecessor below an $n+1$ and $n+2$ block that bound a structure
        isomorphic to $\Z^{\delta_n}$.

    \item Define $z_{n+1}$ analogously to $z_n$, expect it is in $c_{n+1}.$
    \item Say that $z_n\leq x<z_{n+1}$.
    \item Check if $z=z_n$ or not.

    \item In an analogous manner to the case of $\A$, define the automorphism
        orbit of $x$ within its copy of $\Z^{\delta_k}$ or the finite
        block that it is in.

\end{enumerate}

This describes the automorphism orbit of a point and is of complexity less than
$\Sinf{\lambda}$. In order to extend this definition to tuples the only
additional detail needed is the order of the tuple and whether the tuples lie in the same point in $\Q$ or not. This can be done by comparing the various values of "$z$" that emerge for each point in the tuple.

In order to finish the proof, all that remains is to show that
$\SR(\K)=\lambda+2$, or in other words, that it has an orbit that it is not
$\Sinf{\lambda+1}$-definable. Let $c$ be the initial point in $\B$ and $b$ be
some other point in $\B$ in a block of size $1$. We will show that 
$(\K,b)\leq_{\lambda+1}(\K,c)$, so the orbit of $c$ is not $\Sinf{\lambda+1}$
definable in $\K$.

Note that $\K_{>c}\cong \K_{>b}$, so it is sufficient to show that
$\K_{<b}\leq_{\lambda+1} \K_{<c}.$
As $\K_{<c}$ is an initial segment of $\K_{<b}$, for the first turn of the game,
the $\exists$-player can pick whatever points the $\forall$-player picked in
$\K_{<c}$ according to the canonical embedding of $\K_{<c}$ in $\K_{<b}$.
Note that each interval is isomorphic except for the one between the greatest chosen points (call them $a_k$) and $c$ or $b$ respectively. 
$a_k$ may be greater than some finite number of the $c_n$. 
Therefore, up to a finite difference in the choice of sequence $\delta_n$ used
to define $\L_{\lambda,n}$, $(a_k,c)$ is just $\M+\A$ where $\M$ is
$\L_{\lambda,>c}$ or $\L_{\lambda,n,>c}$ depending on which type of order $a_k$
lives in and $(a_k,b)$ is just $\M+\K$.
In short, it is enough to demonstrate that
$$\A\geq_\lambda \K,$$
to show the claim.
To see this, choose $\alpha<\lambda$ and show that
$\A\geq_\alpha \K.$
Take $\delta_n>\alpha$. It is indeed the case that for $k\geq n$,
$\L_{\lambda}\equiv_\alpha \L_{\lambda,k}$. For some point $d$ in a block of
size $1$ between $c_n$ and $c_{n+1}$ we have that
\[\A=\A_{<d}+\sum_{q\in\Q} \A_q~~ \text{ and } ~~ \K=\A_{<d}+\sum_{q\in\Q} \B_q,\]
where $\A_q\equiv_\alpha \L_{\lambda}\equiv_\alpha \B_q$. This proves the claim.
\end{proof}

\begin{thebibliography}{9}
\bibitem{ash2000} C. Ash and J. Knight. \emph{Computable Structures and the Hyperarithmetical Hierarchy}. Studies in Logic and the Foundations of Mathematics144,2000.
\bibitem{AR20} R. Alvir and D. Rossegger. The complexity of Scott sentences of scattered linear orders. \emph{Journal of Symbolic Logic}85(3)(2020),1079--1101.
\bibitem{ash1986} C.J. Ash. Recursive labelling systems and stability of recursive structures in hyperarithmetical degrees. \emph{Transactions of the American Mathematical Society}298(2)(1986),497--514.
\bibitem{montalban2021} A. Montalban. \emph{Computable Structure Theory: Beyond the Arithmetic}. Draft,2021. \texttt{https://math.berkeley.edu/\string~antonio/CSTpart2\_DRAFT.pdf}.
\bibitem{GM23} David Gonzalez and Antonio Montalb'an. \emph{The $\omega$-Vaught's Conjecture}. Native bibliography entry supplied in flattened.bbl.
\bibitem{MonIntermediate} Antonio Montalb'an. Classes of structures with no intermediate isomorphism problems. \emph{Journal of Symbolic Logic}81 (2016),127--150.
\end{thebibliography}
\end{document}
